\section{Related Work}
\label{sec:related_work}

Our work is situated at the intersection of model merging, dynamic routing, and low-data
regularization. In this section, we review the foundational literature in each of these subfields
and articulate how our proposed Task-Space Anchor Regularization (TSAR) differentiates itself.

\subsection{Static Model Merging}
Model merging has emerged as a compelling alternative to joint training or multi-task fine-tuning
by directly fusing the parameters of specialized, pre-trained neural networks. Early, foundational
work introduced \emph{Task Arithmetic} \cite{ilharco2022editing}, which constructs task vectors by
subtracting pre-trained base model weights from task-specific fine-tuned weights and linearly
interpolating them using uniform coefficients. Concurrently, \emph{Model Soups}
\cite{wortsman2022model} demonstrated that averaging the weights of multiple models fine-tuned with
different hyperparameters on the same task significantly improves generalization. To merge models
trained from different random initializations, \emph{Git Re-Basin} \cite{ainsworth2022git} leveraged
permutation matrices to align internal representation channels before parameter averaging, resolving
coordinate-matching conflicts across loss landscapes. To optimize weight-averaging,
\emph{AdaMerging} \cite{yang2023adamerging} parameterized the static merging coefficients and
calibrated them post-hoc using a small validation set. While static merging approaches are
computationally efficient and easy to deploy, they enforce a single, global weight configuration for
all incoming test inputs, which limits their adaptability to heterogeneous or non-stationary input
streams.

\subsection{Dynamic Model Merging and Mixture-of-Experts}
To overcome the inflexibility of static merging, modern frameworks have integrated dynamic routing
mechanisms inspired by Mixture-of-Experts (MoE) architectures \cite{shazeer2017outrageously,
fedus2021switch}. Instead of selecting a fixed set of merging coefficients, dynamic model merging
models utilize a lightweight, input-dependent router to output sample-specific merging coefficients
on the fly \cite{shavit2023quantum, trial5_submission5}. For instance, \emph{QWS-Merge}
\cite{shavit2023quantum} proposes a highly complex, "quantum-inspired" wave-superposition routing
layer that represents coefficients as phase-state transitions, which aims to preserve multi-task
superposition. More recently, the \emph{L3-Router} \cite{trial5_submission5} introduced a
layer-wise, low-dimensional classical routing network that projects inputs onto a low-dimensional
feature subspace to dynamically scale layer-wise merging coefficients. While these dynamic
approaches offer superior theoretical capacity, they rely on a post-hoc calibration phase over small
calibration splits. Prior research has largely ignored the optimization challenges in this
calibration phase, assuming that standard cross-entropy optimization and standard weight decay are
sufficient. In this paper, we expose the catastrophic low-data overfitting vulnerability inherent in
these dynamic routing models and provide a rigorous empirical resolution.

\subsection{Low-Data Regularization and Anchor-Guided Learning}
The challenges of model optimization under extreme data scarcity have been extensively studied in
few-shot and low-data learning literature \cite{snell2017prototypical, loshchilov2017adamw}. Under
low-data constraints, unconstrained optimization parameters are prone to drift and overfit to local
sampling artifacts. To stabilize representations, researchers have proposed prototypical and
anchor-guided alignment constraints. In few-shot classification, \emph{Prototypical Networks}
\cite{snell2017prototypical} compute task centroids (prototypes) in a shared metric space and
classify query samples based on distance, which significantly reduces sample complexity. Similarly,
distance-based representation classifiers \cite{mensink2013distance, bengio2013representation} have
shown that anchoring parameters to stable, pre-trained reference points prevents
representation-space collapse and parameter drift. Drawing inspiration from these techniques, TSAR
adapts the concept of prototype guidance to the parameter-fusion domain. By pre-computing task
feature centroids (anchors) directly from robust, pre-trained expert representations and penalizing
the spatial distance of layer-wise routing weights from these anchors, TSAR provides a highly
stable, lightweight regularization framework that resolves low-data calibration overfitting without
sacrificing routing expressiveness.